\documentclass[11pt]{article}
\usepackage[margin=1in]{geometry}
\usepackage{amsmath,amssymb,amsthm}
\usepackage{graphicx}
\usepackage[hidelinks]{hyperref}
\usepackage{xcolor}

\newtheorem{theorem}{Theorem}
\newtheorem{lemma}{Lemma}
\newcommand{\N}{\mathbb{N}}

\title{A bounded ultrametric counterexample to the\\
right-semicontinuous uniform-convergence claim}
\author{Automatic proof packet for arXiv:2507.05276}
\date{9 August 2026}

\begin{document}
\maketitle

\begin{center}
\fbox{\parbox{0.91\textwidth}{\textbf{Status.}
Candidate counterexample, likely valid, pending human review.  The example
disproves the uniform-on-bounded-sets conclusion of Corollary 3.9 and the
uniform clause of Theorem 3.11 in Khandani's arXiv:2507.05276v1.  It does not
challenge the pointwise-convergence or fixed-point clauses.}}
\end{center}

\section{Source claim}

Khandani's paper \cite{khandani} states in Corollary 3.9 that the bounded-set
uniform-convergence theorem remains valid when upper semicontinuity of the
control function is weakened to upper semicontinuity from the right.  The
corollary depends on Lemma 3.8, which claims that every positive function on a
compact interval that is lower semicontinuous from the right has positive
infimum.  Theorem 3.11 then incorporates the same uniform-convergence claim
under any of three additional fixed-point hypotheses.

The source statements and their proof dependency appear on pages 9--10:

\begin{center}
\includegraphics[width=0.96\textwidth,height=0.88\textheight,keepaspectratio]{figures/open_problem_crop.png}
\end{center}

\clearpage
\begin{center}
\includegraphics[width=0.96\textwidth,height=0.88\textheight,keepaspectratio]{figures/source_theorem_crop.png}
\end{center}

\clearpage
\section{Proof intuition}

Use a bounded ultrametric space consisting of countdown chains of every finite
length.  The point $x_k$ needs exactly $k$ iterates to reach the fixed point
$x_0$.  Distances $r_k=k/(k+1)$ approach the unattained diameter $1$.
One step lowers $r_k$ to $r_{k-1}$, so a control function can lie strictly
below the identity at every distance.  The gaps $r_k-r_{k-1}$ nevertheless
tend to zero as $r_k\uparrow1$.  Right semicontinuity at $1$ imposes no
condition on this left-hand approach.  Consequently every orbit converges
(indeed terminates), while no common iteration count works for the entire
bounded space.

\section{The counterexample}

Let
\[
 r_0=0,\qquad r_k=\frac{k}{k+1}\quad(k\geq1),
\]
and let $X=\{x_k:k\in\N_0\}$.  Define
\[
 d(x_j,x_k)=
 \begin{cases}
 0,&j=k,\\
 r_{\max\{j,k\}},&j\ne k.
 \end{cases}
\]
Define $T:X\to X$ by
\[
 Tx_0=x_0,\qquad Tx_k=x_{k-1}\quad(k\geq1).
\]
Finally, define $\phi:[0,\infty)\to[0,\infty)$ by $\phi(0)=0$ and
\[
 \phi(t)=
 \begin{cases}
 t/2,&0<t<r_1,\\
 1/4,&r_1\leq t<r_2,\\
 r_{k-1},&r_k\leq t<r_{k+1}\quad(k\geq2),\\
 1/2,&t\geq1.
 \end{cases}
\]
For every positive integer $n$, set $\phi_n=\phi$.

\begin{theorem}
The complete metric space $(X,d)$, map $T$, and functions $\phi_n=\phi$
satisfy all hypotheses of Theorem 3.11 in \cite{khandani}.  The map $T$ is
continuous, has a complete graph and unique fixed point $x_0$, and $T^n x$
converges to $x_0$ for every $x\in X$.  However, the convergence is not
uniform on the bounded set $X$.
\end{theorem}

\begin{proof}
First, $d$ is an ultrametric.  Indeed, for three distinct indices, the largest
index determines the two largest pairwise distances, which are equal; cases
with repeated points are immediate.  Every two distinct points are at
distance at least $r_1=1/2$.  Hence every Cauchy sequence is eventually
constant, so $X$ is complete.  Also $\operatorname{diam}X=1$, where the
diameter is a supremum, so $X$ is bounded.  The same separation shows that
$X$ is discrete.  Therefore $T$ is continuous, its graph is closed and,
as a closed subspace of the complete product $X\times X$, its graph is
complete.

The function $\phi$ is positive and strictly below the identity on
$(0,\infty)$.  This is immediate on $(0,r_1)$ and $[1,\infty)$.  On
$[r_1,r_2)$, $\phi=1/4<r_1$, and on
$[r_k,r_{k+1})$ for $k\geq2$, $\phi=r_{k-1}<r_k$.
It is right-continuous at every point except that only the one-sided statement
is needed at the step boundaries; directly from the half-open definition,
the right limit equals the displayed value.  At $t=1$ and for $t>1$ it is
constant from the right.  Thus $\phi$ is upper semicontinuous from the right
at every positive $t$.  Notice that it is not upper semicontinuous at $1$ in
the two-sided sense, since
\[
 \lim_{k\to\infty}\phi(r_k)
 =\lim_{k\to\infty}r_{k-1}=1>\phi(1)=\frac12.
\]
The constant sequence $\phi_n=\phi$ converges uniformly to $\phi$ on every
bounded interval.

It remains to verify the asymptotic contraction inequality.  We have
\[
 T^n x_k=x_{\max\{k-n,0\}}.
\]
For distinct $x_j,x_k$, suppose without loss of generality that $j<k$.
For $n\geq1$, either both iterates are $x_0$, or their distance is
$r_{k-n}\leq r_{k-1}$.  If $k=1$, both iterates already equal $x_0$;
if $k\geq2$, then $\phi(d(x_j,x_k))=\phi(r_k)=r_{k-1}$.  Therefore
\[
 d(T^n x_j,T^n x_k)\leq\phi_n(d(x_j,x_k))
 \qquad(n\geq1).
\]
The case $j=k$ follows from $\phi_n(0)=0$.  Thus all three numbered
hypotheses of source Theorem 3.11 hold.

The only fixed point is $x_0$, and each orbit reaches it after finitely many
steps.  Uniform convergence on $X$ nevertheless fails.  For every $n\geq0$,
\[
 T^n x_{n+4}=x_4,
 \qquad d(T^n x_{n+4},x_0)=r_4=\frac45.
\]
Hence no iteration index makes all orbit points even $3/4$-close to $x_0$.
Equivalently,
\[
 \sup_{x\in X}d(T^n x,x_0)=1
 \qquad\text{for every }n.
\]
This contradicts the asserted bounded-set uniform convergence.  Since $T$
is continuous and has a complete graph, it satisfies both conditions 1 and 3
listed after the conclusion in source Theorem 3.11.
\end{proof}

\section{Why the source proof fails}

Lemma 3.8 in \cite{khandani} is false.  For example, on $[1/2,1]$ let
\[
 f(t)=1-t\quad(1/2\leq t<1),\qquad f(1)=1.
\]
Then $f>0$ and $f$ is lower semicontinuous from the right, but
$\inf f=0$.  The intervals $[t,t+\delta_t)$ used in the source proof are not
open in the usual topology, so compactness does not yield the claimed finite
subcover.  In the constructed dynamical example, the same obstruction is
encoded by
\[
 r_k-\phi(r_k)=r_k-r_{k-1}\longrightarrow0.
\]

Full upper semicontinuity would prevent this behavior on a compact interval:
$t-\phi(t)$ would be positive and lower semicontinuous, hence would have a
positive minimum.  A weaker repair is to assume directly that
$\inf_{t\in[a,b]}(t-\phi(t))>0$ for every $0<a<b<\infty$.  The packet does
not claim that this is the weakest possible repair.

\section{Verification and novelty limits}

The accompanying script uses exact rational arithmetic to check all
ultrametric triples through index $80$, all contraction inequalities for
iterates $1$ through $40$ on that truncation, and the explicit witnesses
$x_{n+4}$ for $0\leq n\leq40$.  These tests check transcription and indexing;
the proof above establishes the infinite statements.

A bounded literature search on 9 August 2026 used the exact source title and
the phrases ``right upper semicontinuity'', ``uniform convergence'',
``Kirk asymptotic contractions'', and ``Leader contraction''.  It found the
source and adjacent positive results but no prior counterexample to the two
identified source claims.  This supports only moderate novelty confidence;
specialist review and a broader citation search remain appropriate.

\section{Limitations and review recommendation}

This counterexample is scoped to bounded-set uniform convergence.  It is
consistent with pointwise convergence, finite-time convergence of every
individual orbit, existence and uniqueness of the fixed point, and the claim
that the particular map is a Leader contraction.  A reviewer should verify
the source's convention $n\geq1$, the values of $\phi$ at $r_k$ and at $1$,
and the exact conclusion of Corollary 3.9/Theorem 3.11.  Subject to those
checks, promotion as a valid counterexample is recommended.

\begin{thebibliography}{9}
\bibitem{khandani}
H. Khandani,
\emph{A Study of Kirk's Asymptotic Contractions via Leader Contractions},
arXiv:2507.05276v1, 2025.
\end{thebibliography}

\end{document}
